\begin{proof}
    First it must be shown that each relation in $E_0$ is sound.
    Let $( X, Y ) \in E_0$.
    Start by noting the following.
    \begin{align*}
        \sem{\scalebox{0.65}{\circimg{tikz-cache/encoder__eqn1_lhs.pdf}}}_H
        &=
        e^{\Log( Z ) \alpha}
        =
        \begin{bsmallmatrix}
            1 & 0 \\
            0 & e^{i\alpha\pi}
        \end{bsmallmatrix}
        =
        Z( \alpha\pi )
        &
        \sem{\scalebox{0.65}{\circimg{tikz-cache/encoder__eqn4_lhs.pdf}}}_H
        &=
        e^{\Log( X ) \alpha} 
        =
        \tfrac{1}{2}
        \begin{bsmallmatrix}
            1 + e^{i\alpha\pi} & 1 - e^{i\alpha\pi} \\
            1 - e^{i\alpha\pi} & 1 + e^{i\alpha\pi}
        \end{bsmallmatrix}
        =
        X( \alpha\pi )
    \end{align*}
    There are several cases to consider.
    \begin{itemize}
    \item Assume that $( X, Y )$ is \HqcHDef. Then the following equation holds.
          \begin{align*}
              \sem{X}_H
              &=
              \sem*{\scalebox{0.65}{\circimg{tikz-cache/hqc__eqn8_lhs.pdf}}}_H
              =
              \tfrac{1}{\sqrt{2}}
              \begin{bsmallmatrix}
                1 & 1 \\
                1 & -1
              \end{bsmallmatrix}
              \\
              &=
              \overbrace{
                \underbrace{
                    \tfrac{1}{2}
                    \begin{bsmallmatrix}
                        1 - i & 1 + i \\
                        1 + i & 1 - i
                    \end{bsmallmatrix}
                }_{X(-\pi/2)}
                \underbrace{
                    \begin{bsmallmatrix}
                        1 & 0 \\
                        0 & e^{i\pi/4}
                    \end{bsmallmatrix}
                }_{Z(\pi/4)}
                \underbrace{
                    \tfrac{1}{2}
                    \begin{bsmallmatrix}
                        1 + i & 1 - i \\
                        1 - i & 1 + i
                    \end{bsmallmatrix}
                }_{X(\pi/2)}
              }^{Y(\pi/4)}
              \underbrace{
                \begin{bsmallmatrix}
                    1 & 0 \\
                    0 & -1
                \end{bsmallmatrix}
              }_{Z}
              \overbrace{
                \underbrace{
                    \tfrac{1}{2}
                    \begin{bsmallmatrix}
                        1 - i & 1 + i \\
                        1 + i & 1 - i
                    \end{bsmallmatrix}
                }_{X(-\pi/2)}
                \underbrace{
                    \begin{bsmallmatrix}
                        1 & 0 \\
                        0 & e^{-i\pi/4}
                    \end{bsmallmatrix}
                }_{Z(-\pi/4)}
                \underbrace{
                    \tfrac{1}{2}
                    \begin{bsmallmatrix}
                        1 + i & 1 - i \\
                        1 - i & 1 + i
                    \end{bsmallmatrix}
                }_{X(\pi/2)}
              }^{Y(-\pi/4)}
              \\
              &=
              \sem*{\scalebox{0.65}{\circimg{tikz-cache/hqc__eqn8_rhs.pdf}}}_H
              =
              \sem{Y}_H
          \end{align*}

    \item Assume that $( X, Y )$ is \HqcTargDef.
          Then the following equation holds.
          \begin{align*}
            \sem{X}_H
            =
            \sem*{\scalebox{0.65}{\circimg{tikz-cache/hqc__eqn10_lhs.pdf}}}_H
            =
            \begin{bsmallmatrix}
                0 & 1 \\
                1 & 0
            \end{bsmallmatrix}
            =
            \tfrac{1}{\sqrt{2}}\begin{bsmallmatrix}
                1 & 1 \\
                1 & -1
            \end{bsmallmatrix}
            \begin{bsmallmatrix}
                1 & 0 \\
                0 & -1
            \end{bsmallmatrix}
            \tfrac{1}{\sqrt{2}}\begin{bsmallmatrix}
                1 & 1 \\
                1 & -1
            \end{bsmallmatrix}
            =
            \sem*{\scalebox{0.65}{\circimg{tikz-cache/hqc__eqn10_rhs.pdf}}}_H
            =
            \sem{Y}_H
          \end{align*}

    \item Assume that $( X, Y )$ is \HqcOPhaseDef.
          Then the following equation holds.
          \begin{align*}
            \sem{X}_H
            =
            \sem*{\scalebox{0.65}{\circimg{tikz-cache/hqc__eqn11_lhs.pdf}}}_H
            =
            \begin{bsmallmatrix}
                -1 & 0 \\
                0 & 1
            \end{bsmallmatrix}
            =
            \begin{bsmallmatrix}
                0 & 1 \\
                1 & 0
            \end{bsmallmatrix}
            \begin{bsmallmatrix}
                1 & 0 \\
                0 & -1
            \end{bsmallmatrix}
            \begin{bsmallmatrix}
                0 & 1 \\
                1 & 0
            \end{bsmallmatrix}
            =
            \sem*{\scalebox{0.65}{\circimg{tikz-cache/hqc__eqn11_rhs.pdf}}}_H
            =
            \sem{Y}_H
          \end{align*}

    \item Assume that $( X, Y )$ is \HqcTargMDef.
          Then the following equation holds.
          \begin{align*}
            \sem{X}_H
            =
            \sem*{\scalebox{0.65}{\circimg{tikz-cache/hqc__eqn12_lhs.pdf}}}_H
            =
            \begin{bsmallmatrix}
                0 & -1 \\
                -1 & 0
            \end{bsmallmatrix}
            =
            \begin{bsmallmatrix}
                1 & 0 \\
                0 & -1
            \end{bsmallmatrix}
            \begin{bsmallmatrix}
                0 & 1 \\
                1 & 0
            \end{bsmallmatrix}
            \begin{bsmallmatrix}
                1 & 0 \\
                0 & -1
            \end{bsmallmatrix}
            =
            \sem*{\scalebox{0.65}{\circimg{tikz-cache/hqc__eqn12_rhs.pdf}}}_H
            =
            \sem{Y}_H
          \end{align*}

    \item Assume that $( X, Y )$ is \HqcSwapDef.
          Then the following equation holds.
          \begin{align*}
            \sem{X}_H
            =
            \sem*{\scalebox{0.65}{\circimg{tikz-cache/proofs__decode_q3_lhs.pdf}}}_H
            &=
            \begin{bsmallmatrix}
                1 & 0 & 0 & 0 \\
                0 & 0 & 1 & 0 \\
                0 & 1 & 0 & 0 \\
                0 & 0 & 0 & 1
            \end{bsmallmatrix}
            \\
            &=
            \begin{bsmallmatrix}
                1 & 0 & 0 & 0 \\
                0 & 1 & 0 & 0 \\
                0 & 0 & 0 & 1 \\
                0 & 0 & 1 & 0
            \end{bsmallmatrix}
            \begin{bsmallmatrix}
                1 & 0 & 0 & 0 \\
                0 & 0 & 0 & 1 \\
                0 & 0 & 1 & 0 \\
                0 & 1 & 0 & 0
            \end{bsmallmatrix}
            \begin{bsmallmatrix}
                1 & 0 & 0 & 0 \\
                0 & 1 & 0 & 0 \\
                0 & 0 & 0 & 1 \\
                0 & 0 & 1 & 0
            \end{bsmallmatrix}
            \\
            &=
            \sem{Z \odot X}_H \sem{X \odot Z}_H \sem{Z \odot X}_H
            \\
            &=
            \sem*{\scalebox{0.65}{\circimg{tikz-cache/proofs__decode_q3_mid.pdf}}}_H
            =
            \sem{Y}_H
          \end{align*}

    \item Assume that $( X, Y )$ is an instance of the relation \HqcEuler.
          Then let $\theta_1, \theta_2 \in \mathbb{R}$ and define $( \phi_0\pi, \phi_1\pi, \phi_2\pi, \phi_3\pi ) = \mathsf{Euler}( \theta_1\pi, \theta_2\pi )$ as in an arbitrary instance of \HqcEuler.
          Then by the soundness of Q(5), the equation $H Z( \theta_1\pi ) H Z( \theta_2\pi ) H = e^{i\phi_0\pi} Z( \phi_1\pi ) H Z( \phi_2\pi ) H Z( \phi_3\pi )$ holds.
          Then the following equation holds.
          \begin{align*}
            \sem{X}_H
            =
            \sem*{\scalebox{0.65}{\circimg{tikz-cache/hqc__euler_lhs_1line.pdf}}}_H
            &=
            H Z( \theta_2\pi ) X( \theta_1\pi )
            \\
            &=
            H Z( \theta_2\pi ) H Z( \theta_1\pi ) H
            \\
            &=
            e^{i\phi_0\pi} Z( \phi_1\pi ) H Z( \phi_2\pi ) H Z( \phi_3\pi )
            \\
            &=
            e^{i\phi_0\pi} Z( \phi_1\pi ) X( \phi_2\pi ) Z( \phi_3\pi )
            \\
            &=
            \sem*{\scalebox{0.65}{\circimg{tikz-cache/hqc__euler_rhs.pdf}}}_H
            =
            \sem{Y}_H
          \end{align*}

    \item Assume that $( X, Y )$ is \HqcPOrd.
          Then the following equation holds.
          \begin{equation*}
            \sem{X}_H
            =
            \sem*{\scalebox{0.65}{\circimg{tikz-cache/hqc__eqn7_lhs.pdf}}}_H
            =
            \exp( 2 \Log\sem{\scalebox{0.65}{\circimg{tikz-cache/dagger__term8.pdf}}}_H )
            =
            \exp( 2 \Log(-1) )
            =
            \exp( 2 (i\pi) )
            =
            \mathbb{I}_1
            =
            \sem{1_0}_H
            =
            \sem{Y}_H
          \end{equation*}

    \item Assume that $( X, Y )$ is an instance of \HqcCtrlId.
          Then the following holds by \cref{Cor:GenCtrlFunc}.
          \begin{equation*}
            \sem{X}_H
            =
            \sem*{\scalebox{0.65}{\circimg{tikz-cache/hqc__eqn1_lhs.pdf}}}_H
            =
            \sem{U}_H \odot \sem{1_m}_H
            =
            \mathbb{I}_{2^{n+m}}
            =
            \mathbb{I}_{2^n} \otimes \mathbb{I}_{2^m}
            =
            \sem*{\scalebox{0.65}{\circimg{tikz-cache/hqc__eqn1_rhs.pdf}}}_H
            =
            \sem{Y}_H
          \end{equation*}
   
    \item Assume that $( X, Y )$ is an instance of \HqcCtrlComp.
          Since $Z$ is unitary with eigenvalues $\{ -1, 1 \}$, then $Z$ is Hermitian.
          Then the following equation holds.
          \begin{align*}
            \sem{X}_H
            =
            \sem*{\scalebox{0.65}{\circimg{tikz-cache/hqc__eqn2_lhs.pdf}}}_H
            &=
            Z \odot \sem{V \circ U}_H
            \\
            &=
            Z \odot ( \sem{V}_H \circ \sem{U}_H )
            \\
            &=
            ( Z \odot \sem{V}_H ) \circ ( Z \odot \sem{U} )_H
            &
            ( \text{by Cor.~\cref{Cor:GenCtrlFunc}} )
            \\
            &=
            \sem*{\scalebox{0.65}{\circimg{tikz-cache/hqc__eqn2_rhs.pdf}}}_H
            =
            \sem{Y}_H
          \end{align*}
          
    \item Assume that $( X, Y )$ is an instance of \HqcCtrlSwap.
          Since $\sem{-}_H: \CatHQC \to \CatUnitary$ is symmetric monoidal, then $\sem*{\scalebox{0.5}{\circimg{tikz-cache/soundness__swap.pdf}}}_H$ is the monoidal symmetry $\sigma_{2^n,2^m}: \mathbb{C}^{2^n} \otimes \mathbb{C}^{2^m} \cong \mathbb{C}^{2^m} \otimes \mathbb{C}^{2^n}$.
          Then $\sigma_{2^n,2^m}$ is the self-adjoint unitary operator whose action is $\sigma_{2^n,2^m} \cdot ( N \otimes M ) \cdot \sigma_{2^n,2^m} = M \otimes N$ for all $M \in \CatUnitary( 2^n, 2^n )$ and $N \in \CatUnitary( 2^m, 2^m )$.
          Then the following equation holds.
          \begin{align*}
            \sem{X}
            =
            \sem*{\scalebox{0.65}{\circimg{tikz-cache/hqc__eqn3_lhs.pdf}}}_H
            &=
            \exp( \Log\sem{U}_H \otimes \Log\sem{V}_H / (i\pi) )
            \\
            &=
            \exp( \sigma_{2^n,2^m} \cdot ( \Log\sem{V}_H \otimes \Log\sem{U}_H / (i\pi) ) \cdot \sigma_{2^n,2^m} )
            \\
            &=
            \sigma_{2^n,2^m} \cdot \exp( \Log\sem{V}_H \otimes \Log\sem{U}_H / (i\pi) ) \cdot \sigma_{2^n,2^m}
            &
            ( \text{by~\cref{Eqn:UnitaryConj}} )
            \\
            &=
            \sem*{\scalebox{0.65}{\circimg{tikz-cache/hqc__eqn3_rhs.pdf}}}_H
            =
            \sem{Y}_H
          \end{align*}

    \item Assume that $( X, Y )$ is an instance of \HqcFreePhase.
          Then the following holds by \cref{Thm:PowAsCtrl}.
          \begin{align*}
            \sem{X}_H
            =
            \sem*{\scalebox{0.65}{\circimg{tikz-cache/hqc__eqn18_lhs.pdf}}}_H
            =
            Z \odot \sem*{\scalebox{0.65}{\circimg{tikz-cache/encoder__eqn3_lhs.pdf}}}
            =
            Z \odot e^{i\alpha\pi}
            =
            Z^\alpha
            =
            \sem*{\scalebox{0.65}{\circimg{tikz-cache/encoder__eqn1_lhs.pdf}}}_H
            =
            \sem{Y}_H
          \end{align*}

    \item Assume that $( X, Y )$ is an instance of \HqcCtrlDag.
          Then the following equation holds.
          \begin{align*}
            \sem{X}_H
            =
            \sem*{\scalebox{0.65}{\circimg{tikz-cache/hqc__eqn14_lhs.pdf}}}_H
            &=
            \sem{V}_H \odot \sem{P^\dagger \circ U \circ P}
            \\
            &=
            \sem{V}_H \odot ( \sem{P^\dagger}_H \circ \sem{U}_H \circ \sem{P}_H )
            \\
            &=
            \sem{V}_H \odot ( \sem{P}_H^\dagger \circ \sem{U}_H \circ \sem{P}_H )
            &
            ( \text{by Thm.~\ref{Thm:DagSem}} )
            \\
            &=
            ( \mathbb{I}_{2^m} \otimes \sem{P}_H^\dagger ) ( \sem{V} \odot \sem{U} ) ( \mathbb{I}_{2^m} \otimes \sem{P}_H )
            &
            ( \text{by Cor.~\ref{Cor:GenCtrlFunc}} )
            \\
            &=
            ( \mathbb{I}_{2^m} \otimes \sem{P^\dagger}_H ) ( \sem{V} \odot \sem{U} ) ( \mathbb{I}_{2^m} \otimes \sem{P}_H )
            &
            ( \text{by Thm.~\ref{Thm:DagSem}} )
            \\
            &=
            \sem*{\scalebox{0.65}{\circimg{tikz-cache/hqc__eqn14_rhs.pdf}}}_H
            =
            \sem{Y}_H
          \end{align*}

    \item Assume that $( X, Y )$ is an instance of \HqcCtrlPair.
          First, note that for each $M \in \CatUnitary( 2^n, 2^n )$, the matrix $M \otimes \ket{0} \bra{0}$ commutes with the matrix $M \otimes \ket{1} \bra{1}$ since $\braket{0|1} = 0 = \braket{1|0}$.
          Then the following equation holds where $A = \Log( -Z ) = i \ket{0} \bra{0} \pi$ and $B = \Log( Z ) = i \ket{1} \bra{1} \pi$.
          \begin{align*}
            \sem{X}_H
            =
            \sem*{\scalebox{0.65}{\circimg{tikz-cache/hqc__eqn15_lhs.pdf}}}_H
            %\\
            &=
            \exp( \Log\sem{U}_H \otimes A / (i\pi) ) \exp( \Log\sem{U}_H \otimes B / (i\pi) )
            \\
            &=
            \exp( \Log\sem{U}_H \otimes \ket{0}\bra{0} ) \exp( \Log\sem{U}_H \otimes \ket{1}\bra{1} )
            &
            ( \text{by~\cref{Eqn:SimulDiag}} )
            \\
            &=
            \exp\big( ( \Log\sem{U}_H \otimes \ket{0}\bra{0} ) + ( \Log\sem{U}_H \otimes \ket{1}\bra{1} ) \big) 
            \\
            &=
            \exp( \Log\sem{U}_H \otimes (\ket{0}\bra{0} + \ket{1}\bra{1}) )
            \\
            &=
            \exp( \Log\sem{U}_H \otimes \mathbb{I}_2 )
            \\
            &=
            \exp( \Log\sem{U}_H ) \otimes \mathbb{I}_2
            \\
            &=
            \sem{U}_H \otimes \mathbb{I}_2
            \\
            &=
            \sem*{\scalebox{0.65}{\circimg{tikz-cache/hqc__eqn15_rhs.pdf}}}_H
            =
            \sem{Y}_H
          \end{align*}

    \item Assume that $( X, Y )$ is an instance of \HqcCtrlExpand.
          As in \cref{App:DiagGate}, let $\ket{x} = \bigotimes_{j=0}^{n-1} \ket{x_j}$.
          Moreover, let $\Lambda$ be the prefix of a diagonal form for which the term $\lambda( x, \rho )$ does not appear.
          This means that the eigenvalue of $\ket{x} \in \sem{\Lambda}_H$ is $1$, and consequently $\sem{\Lambda}_H$ has the form $\sem{\Lambda}_H = \exp\left( \sum_{y \in \mathcal{B}} i \theta_y \ket{y} \bra{y} \right)$ where $\mathcal{B} = \{ 0, 1 \}^n \setminus \{ x \}$ and $\theta_y \in ( -\pi, \pi ]$ for each $y \in \mathcal{B}$ (see~\cref{Cor:Diag}).
          Since $x \not \in \mathcal{B}$, then for each $y \in \mathcal{B}$, $\braket{x|y} = 0 = \braket{y|x}$ and consequently the matrix $\ket{x}\bra{x}$ commutes with $\ket{y}\bra{y}$.
          Then the following equation holds for $M = \Log\sem{\Gamma}_H / (i\pi)$.
          \begin{align*}
            \sem{X}_H
            &=
            \sem*{\scalebox{0.65}{\circimg{tikz-cache/hqc__eqn22_lhs.pdf}}}_H
            \\
            &=
            \exp\left( \Log\sem{\Lambda \circ \lambda(x,\rho)}_H \otimes M \right)
            \\
            &=
            \exp\left( \Log\left( \sem{\Lambda}_H \sem{\lambda(x,\rho)}_H  \right) \otimes M \right)
            \\
            &=
            \exp\left( \Log\left( \exp\left( \sum_{y \in \mathcal{B}} i \theta_y \ket{y}\bra{y} \right) \exp( i\rho\ket{x}\bra{x} ) \right) \otimes M \right)
            &
            ( \text{by Thm.~\ref{Thm:LambdaSem}} )
            \\
            &=
            \exp\left( \Log\left( \exp\left( \sum_{y \in \mathcal{B}} \left( i \theta_y \ket{y}\bra{y} \right) + i\rho \ket{x}\bra{x} \right) \right) \otimes M \right)
            &
            ( \text{by \cref{Eqn:SimulDiag}} )
            \\
            &=
            \exp\left( \Log\left( \sum_{y \in \mathcal{B}} \left( e^{i \theta_y} \ket{y}\bra{y} \right) + e^{i \rho} \ket{x}\bra{x} \right) \otimes M \right)
            \\
            &=
            \exp\left( \left( \sum_{y \in \mathcal{B}} \left( \Log\left( e^{i \theta_y} \right) \ket{y}\bra{y} \right) + \Log\left( e^{i \rho} \right) \ket{x}\bra{x} \right) \otimes M \right)
            \\
            &=
            \exp\left( \left( \left( \sum_{y \in \mathcal{B}} i \theta_y \ket{y}\bra{y} \right) + i \rho \ket{x}\bra{x} \right) \otimes M \right)
            \\
            &=
            \exp\left( \left( \Log\sem{\Lambda}_H + \Log\sem{\lambda(x,\rho)}_H \right) \otimes M \right)
            \\
            &=
            \exp \left( \left( \Log\sem{\Lambda}_H \otimes M \right) + \left( \Log\sem{\lambda(x,\rho)}_H \otimes M \right) \right)
            \\
            &=
            \exp \left( \Log\sem{\Lambda}_H \otimes M \right) \exp \left( \Log\sem{\lambda(x,\rho)}_H \otimes M \right)
            &
            ( \text{by \cref{Eqn:SimulDiag}} )
            \\
            &=
            \sem*{\scalebox{0.65}{\circimg{tikz-cache/hqc__eqn22_rhs.pdf}}}_H
            =
            \sem{Y}_H
          \end{align*}
          
    \item Assume that $( X, Y )$ is an instance of \HqcCtrlMerge.
          Then let $\rho \in (-\pi, \pi]$ and $\gamma \in (-\pi, \pi]$ as in an arbitrary instance of \HqcCtrlMerge.
          Then the following equation holds.
          \begin{align*}
            \sem{X}_H
            =
            \sem*{\scalebox{0.65}{\circimg{tikz-cache/hqc__eqn23_lhs.pdf}}}_H
            &=
            \exp( \Log\sem{\lambda( x, \rho )}_H \otimes \Log\sem{\lambda( y, \gamma )}_H / ( i\pi ) )
            \\
            &=
            \exp\left( (i \rho\ket{x}\bra{x}) \otimes (i\gamma\ket{y}\bra{y}) / ( i\pi ) \right)
            &
            ( \text{by Thm.~\ref{Thm:LambdaSem}} )
            \\
            &=
            \exp\left( \left( i^2 \rho\gamma \ket{xy} \bra{xy} \right) / ( i\pi ) \right)
            \\
            &=
            \exp( i ( \rho \gamma / \pi ) \ket{xy}\bra{xy} )
            \\
            &=
            \sem*{\scalebox{0.65}{\circimg{tikz-cache/hqc__eqn23_rhs.pdf}}}_H
            =
            \sem{Y}_H
          \end{align*}

    \item Assume that $( X, Y )$ is an instance of \HqcPowZero.
          Then the following equation holds.
          \begin{align*}
            \sem{X}_H
            =
            \sem*{\scalebox{0.65}{\circimg{tikz-cache/hqc__eqn9_lhs.pdf}}}_H
            =
            \exp( 0 \Log\sem{U}_H )
            =
            \exp( \mathbb{O}_{2^n} )
            =
            \mathbb{I}_{2^n}
            =
            \sem*{\scalebox{0.65}{\circimg{tikz-cache/hqc__eqn9_rhs.pdf}}}_H
            =
            \sem{Y}_H
          \end{align*}
    
    \item Assume that $( X, Y )$ is an instance of \HqcPowId.
          Then the following equation holds.
          \begin{align*}
            \sem{X}_H
            =
            \sem*{\scalebox{0.65}{\circimg{tikz-cache/hqc__eqn4_lhs.pdf}}}_H
            =
            \exp( 1 \Log\sem{U}_H )
            =
            \exp( \Log\sem{U}_H )
            =
            U
            =
            \sem*{\scalebox{0.65}{\circimg{tikz-cache/hqc__eqn4_rhs.pdf}}}_H
            =
            \sem{Y}_H
          \end{align*}

    \item Assume that $( X, Y )$ is an instance of \HqcPowAdd.
          Then the following equation holds.
          \begin{align*}
            \sem{X}_H
            =
            \sem*{\scalebox{0.65}{\circimg{tikz-cache/hqc__eqn5_lhs.pdf}}}_H
            &=
            \exp( \beta \Log\sem{U}_H ) \exp( \alpha \Log\sem{U}_H )
            \\
            &=
            \exp( (\alpha + \beta) \Log\sem{U}_H )
            &
            ( \text{by~\cref{Eqn:SimulDiag}} )
            \\
            &=
            \sem*{\scalebox{0.65}{\circimg{tikz-cache/hqc__eqn5_rhs.pdf}}}_H
            =
            \sem{Y}_H
          \end{align*}

    \item Assume that $( X, Y )$ is an instance of \HqcPowIdem.
          Then the following equation holds.
          \begin{align*}
            \sem{X}_H
            =
            \sem*{\scalebox{0.65}{\circimg{tikz-cache/hqc__eqn24_lhs.pdf}}}_H
            =
            \exp( \alpha \Log( \mathbb{I}_{2^n} ) )
            =
            \exp( \alpha \mathbb{O}_{2^n} )
            =
            \exp( \mathbb{O}_{2^n} )
            =
            \mathbb{I}_{2^n}
            =
            \sem*{\scalebox{0.65}{\circimg{tikz-cache/hqc__eqn24_rhs.pdf}}}_H
            =
            \sem{Y}_H
          \end{align*}

    \item Assume that $( X, Y )$ is an instance of \HqcPowDag.
          Then the following equation holds.
          \begin{align*}
            \sem{X}_H
            =
            \sem*{\scalebox{0.65}{\circimg{tikz-cache/hqc__eqn19_lhs.pdf}}}_H
            &=
            \sem{P^\dagger}_H \exp( \Log\sem{U}_H \alpha ) \sem{P}_H
            \\
            &=
            \sem{P}_H^\dagger \exp( \Log\sem{U}_H \alpha ) \sem{P}_H
            &
            ( \text{by Thm.~\ref{Thm:DagSem}} )
            \\
            &=
            \exp\left( \sem{P}_H^\dagger \Log\sem{U}_H \sem{P}_H \alpha \right) 
            &
            ( \text{by~\cref{Eqn:UnitaryConj}} )
            \\
            &=
            \exp\left( \Log \left( \sem{P}_H^\dagger \sem{U}_H \sem{P}_H \right) \alpha \right)
            &
            ( \text{by~\cref{Eqn:UnitaryConj}} )
            \\
            &=
            \sem*{\scalebox{0.65}{\circimg{tikz-cache/hqc__eqn19_rhs.pdf}}}_H
            =
            \sem{Y}_H
          \end{align*}

    \item Assume that $( X, Y )$ is an instance of \HqcPowExpand.
          As in \cref{App:DiagGate}, let $\ket{x} = \bigotimes_{j=0}^{n-1} \ket{x_j}$.
          Moreover, let $\Lambda$ be the prefix of a diagonal form for which the term $\lambda( x, \rho )$ does not appear.
          This means that the eigenvalue of $\ket{x} \in \sem{\Lambda}_H$ is $1$, and consequently $\sem{\Lambda}_H$ has the form $\sem{\Lambda}_H = \exp\left( \sum_{y \in \mathcal{B}} i \theta_y \ket{y} \bra{y} \right)$ where $\mathcal{B} = \{ 0, 1 \}^n \setminus \{ x \}$ and $\theta_y \in ( -\pi, \pi ]$ for each $y \in \mathcal{B}$ (see~\cref{Cor:Diag}).
          Since $x \not \in \mathcal{B}$, then for each $y \in \mathcal{B}$, $\braket{x|y} = 0 = \braket{y|x}$ and consequently the matrix $\ket{x}\bra{x}$ commutes with $\ket{y}\bra{y}$.
          Then the following equation holds.
          \begin{align*}
            \sem{X}_H
            &=
            \sem*{\scalebox{0.65}{\circimg{tikz-cache/hqc__eqn21_lhs.pdf}}}_H
            \\
            &=
            \exp\left( \Log\sem{\Lambda \circ \lambda(x,\rho)}_H \alpha \right)
            \\
            &=
            \exp\left( \Log\left( \sem{\Lambda}_H \sem{\lambda(x,\rho)}_H  \right) \alpha \right)
            \\
            &=
            \exp\left( \Log\left( \exp\left( \sum_{y \in \mathcal{B}} i \theta_y \ket{y}\bra{y} \right) \exp( i\rho\ket{x}\bra{x} ) \right) \alpha \right)
            &
            ( \text{by Thm.~\ref{Thm:LambdaSem}} )
            \\
            &=
            \exp\left( \Log\left( \exp\left( \sum_{y \in \mathcal{B}} \left( i \theta_y \ket{y}\bra{y} \right) + i\rho \ket{x}\bra{x} \right) \right) \alpha \right)
            &
            ( \text{by \cref{Eqn:SimulDiag}} )
            \\
            &=
            \exp\left( \Log\left( \sum_{y \in \mathcal{B}} \left( e^{i \theta_y} \ket{y}\bra{y} \right) + e^{i \rho} \ket{x}\bra{x} \right) \alpha \right)
            \\
            &=
            \exp\left( \left( \sum_{y \in \mathcal{B}} \left( \Log\left( e^{i \theta_y} \right) \ket{y}\bra{y} \right) + \Log\left( e^{i \rho} \right) \ket{x}\bra{x} \right) \alpha \right)
            \\
            &=
            \exp\left( \sum_{y \in \mathcal{B}} \left( i \theta_y \alpha \ket{y}\bra{y} \right) + i \alpha \rho \ket{x}\bra{x} \right)
            \\
            &=
            \exp \left( \sum_{y \in \mathcal{B}} i \theta_y \alpha \ket{y}\bra{y} \right) \exp \Big( i \alpha \rho \ket{x}\bra{x} \Big)
            &
            ( \text{by \cref{Eqn:SimulDiag}} )
            \\
            &=
            \exp\left( \Log \sem{\Lambda}_H \alpha \right) \cdot \sem{\lambda(x,\alpha\rho)}_H
            \\
            &=
            \sem*{\scalebox{0.65}{\circimg{tikz-cache/hqc__eqn21_rhs.pdf}}}_H
            =
            \sem{Y}_H
          \end{align*}
    \end{itemize}
    Since $( X, Y )$ was arbitrary, then $\sem{-}_H$ is sound with respect to $E_0$.
    It then follows by induction on $E_n$ that $E$ is sound with respect to $\sem{-}_H$.
    \begin{itemize}
    \item \textbf{Base Case}.
             It follows by the proceeding analysis that $E_0$ is sound with respect to $\sem{-}_H$.
    \item \textbf{Inductive Hypothesis}.
             For some $n \in \mathbb{N}$, $E_n$ is sound with respect to $\sem{-}_H$.
    \item \textbf{Inductive Step}.
          Assume that for some $n \in \mathbb{N}$, that $E_n$ is sound with respect to $\sem{-}_H$.
          Then let $( X, Y ) \in E_{n+1}$.
          There are three cases to consider.
          \begin{enumerate}
          \item Assume that $( X, Y ) \in E_n$.
                Then $\sem{X}_H = \sem{Y}_H$ by the inductive hypothesis.
          \item Assume that there exists some $U \in \Mor( \Core )$ and $V \approx_{E_n} W$ such that $X = U \odot V$ and $Y = U \odot W$.
                Then by the inductive hypothesis, $\sem{V}_H = \sem{W}_H$.
                It follows that $\sem{X}_H = \exp( \Log\sem{U}_H \otimes \sem{V}_H / (i\pi) ) = \exp( \Log\sem{U}_H \otimes \sem{W}_H / (i\pi) ) = \sem{Y}_H$.
          \item Assume that there exists some exponent $\alpha \in \mathbb{R}$ and $V \approx_{E_n} W$ such that $X = V \exparr \alpha$ and $Y = W \exparr \alpha$.
                Then by the inductive hypothesis, $\sem{V}_H = \sem{W}_H$.
                It follows that $\sem{X}_H = \exp( \alpha \Log\sem{U}_H ) = \exp( \alpha \Log\sem{W}_H ) = \sem{Y}_H$.
          \end{enumerate}
          In each case, $\sem{X}_H = \sem{Y}_H$.
          Since $( X, Y )$ was arbitrary, then $E_{n+1}$ is sound with respect to $\sem{-}_H$.
          Then the inductive step holds.
    \end{itemize}
    Then by the principle of inductive, $E_n$ is sound with respect to $\sem{-}_H$ for each $n \in \mathbb{N}$.
    Assume that $( X, Y ) \in E$.
    Then there exists some $n \in \mathbb{N}$ such that $( X, Y ) \in E_n$.
    Then $\sem{X}_H = \sem{Y}_H$.
    Since $( X, Y )$ was arbitrary, then $E$ is sound with respect to $\sem{-}_H$.
\end{proof}